\chapter{Des couleurs et des histoires}
\label{chap:des-couleurs-et-des-histoires}

Rielle voyageait avec la caravane depuis quelques jours lorsque Kamatsu
commença à l'observer. Elle possédait plusieurs carnets, quelques pinceaux et
une petite boîte de pigments qu'elle sortait lors des haltes. Elle peignait des
paysages, des visages, parfois une fleur ou un vieux pont aperçu quelques
heures auparavant.

Kamatsu commença par regarder de loin, puis d'un peu moins loin, jusqu'au jour
où Rielle posa son pinceau et tourna la tête vers lui.

— Tu sais que je t'ai vu ?

— Je ne voulais pas vous déranger.

Elle sourit.

— Tu ne me déranges pas. Allez, approche.

Kamatsu vint immédiatement s'asseoir à côté d'elle et observa la peinture
représentant un paysage traversé plus tôt dans la journée. Après quelques
secondes, il regarda le tableau, puis l'horizon, puis de nouveau le tableau.

— Ce n'était pas exactement comme ça.

— Non.

— Vous vous êtes trompée ?

— Non plus.

Devant son incompréhension, Rielle sourit.

— Je ne peins pas ce qu'il y avait. Je peins ce que j'en ai retenu.

Avant que Kamatsu puisse approfondir cette étrange idée, Bartiméus vint
s'installer exactement entre Rielle et son œuvre. Elle le regarda, le chat lui
rendit son regard, puis elle prit une nouvelle feuille.

— Très bien. Puisque tu insistes.

\scenebreak

Faire poser Bartiméus se révéla évidemment impossible. Le chat resta immobile
moins d'une minute avant de partir, de revenir, de changer de position,
d'entreprendre sa toilette puis de disparaître de nouveau. Le lendemain,
Kamatsu fut donc très surpris de voir Rielle continuer son portrait alors que
le modèle n'était même pas présent.

— Comment faites-vous ?

Rielle lui demanda de fermer les yeux et d'imaginer Bartiméus.

\illustration
  {1400.png}
  {Rielle et Bartiméus}
  {rielle-et-bartimeus}

— Quelle couleur ont ses yeux ?

Kamatsu répondit.

— Et sa queue ? Comment est-ce qu'il la tient quand il est content ?

Il réfléchit encore.

— Et ses oreilles quand il est contrarié ?

Peu à peu, Kamatsu réalisa qu'il pouvait parfaitement voir le chat sans que
celui-ci soit devant lui. Il connaissait la couleur de ses yeux, la position
de ses oreilles, la manière dont sa queue bougeait et quantité de petits
détails auxquels il n'avait jamais vraiment pensé auparavant.

— Regarder, c'est une chose, lui expliqua Rielle. Observer, c'en est une autre.
Et se souvenir de ce qu'on a ressenti en regardant... encore une autre.

Kamatsu demanda à essayer. Son premier Bartiméus ressemblait assez peu à
Bartiméus. Le deuxième aussi. Le troisième ressemblait au moins davantage à un
chat, ce que Rielle considéra comme un progrès.

Kael tenta lui aussi l'expérience une fois. Après plusieurs minutes d'efforts,
Kamatsu examina son dessin.

— C'est Bartiméus ?

— Oui.

— D'accord.

— Pourquoi ?

— Pour rien.

Rielle devint soudain extrêmement intéressée par sa palette et Kael conclut que
la peinture ne serait probablement pas son domaine.

\scenebreak

Kamatsu, en revanche, continua. Chaque fois que Rielle sortait ses pinceaux ou
son fusain, il venait observer ce qu'elle faisait et essayait à son tour. Elle
lui montra comment construire une forme avant de chercher les détails, comment
quelques traits pouvaient suffire à suggérer un mouvement et comment les
couleurs changeaient selon la lumière.

Mais elle revenait toujours à la même chose : avant de vouloir représenter
quelque chose, il fallait apprendre à l'observer.

Kamatsu commença ainsi à regarder différemment ce qui l'entourait. Un arbre
n'était plus seulement un arbre lorsqu'il essayait de le dessiner. Il fallait
voir la forme de son tronc, la direction de ses branches, la manière dont la
lumière traversait ses feuilles. Il en allait de même pour un visage, un
chariot ou un paysage aperçu depuis la route.

Ses dessins restaient maladroits, mais ils s'amélioraient.

Avant son départ, Rielle offrit à Kamatsu le portrait terminé de Bartiméus
ainsi que quelques feuilles et un morceau de fusain.

— Pour que je continue ? demanda-t-il.

— Non. Pour que tu puisses continuer si tu en as envie.

Kamatsu regarda le portrait, puis le matériel qu'elle venait de lui donner.

— Merci.

Il conserva longtemps le portrait.